% ============================================================
\chapter{Dimensional Analysis and Scaling}
\label{ch:dimensional}
% ============================================================

\section{Introduction}

In 1883, Osborne Reynolds was studying the flow of water through pipes at the University of Manchester. He noticed that the transition from smooth, laminar flow to chaotic, turbulent flow did not depend separately on the pipe diameter, the flow speed, or the fluid viscosity---it depended on a single dimensionless combination of all three, now called the \emph{Reynolds number}. This was not a coincidence. It was a manifestation of a deep principle: the laws of physics cannot depend on the arbitrary units we choose to measure them in.

Dimensional analysis turns this seemingly obvious observation into a remarkably powerful tool. The Buckingham $\Pi$ theorem, formalised by Edgar Buckingham in 1914, states that any physically meaningful relationship among $n$ variables involving $k$ independent dimensions can be rewritten as a relationship among $n - k$ dimensionless groups. The consequences are far-reaching: models are simplified, the number of experiments needed is drastically reduced, and scaling laws emerge naturally. From wind tunnel tests of aircraft to the energy of nuclear explosions (famously estimated by G.~I.~Taylor from a single photograph), dimensional analysis reveals structure hidden in plain sight.

% ------------------------------------------------------------
\section{Dimensions and Units}
% ------------------------------------------------------------

\begin{definition}[Physical dimension]
Every physical quantity possesses a \textbf{dimension} expressed in
terms of fundamental dimensions.  In the SI system, the base
dimensions are:
\begin{center}
\begin{tabular}{ccc}
\toprule
Quantity & Dimensional symbol & SI unit \\
\midrule
Length & $L$ & m \\
Mass & $M$ & kg \\
Time & $T$ & s \\
Temperature & $\Theta$ & K \\
Amount of substance & $N$ & mol \\
Electric current & $I$ & A \\
\bottomrule
\end{tabular}
\end{center}
\end{definition}

\begin{definition}[Dimensional homogeneity]
A physical equation is \textbf{dimensionally homogeneous} if every
term has the same dimension.  This is a necessary (but not sufficient)
condition for validity.
\end{definition}

\begin{example}
Velocity $v$ has dimension $[v]=LT^{-1}$.  Kinetic energy
$E_k = \tfrac{1}{2}mv^2$ has dimension $[E_k]=ML^2T^{-2}$, which is
indeed the dimension of energy.
\end{example}

\begin{remark}
Arguments of transcendental functions ($\exp$, $\ln$, $\sin$, etc.)
must be dimensionless.  Thus in $e^{-t/\tau}$, the ratio $t/\tau$ is
dimensionless.
\end{remark}

% ------------------------------------------------------------
\section{The Buckingham $\Pi$ Theorem}
\label{sec:buckingham}
% ------------------------------------------------------------

\begin{theorem}[Buckingham $\Pi$]
Let a physical relationship involve $n$ quantities $q_1,\ldots,q_n$
whose dimensions are expressed using $k$ independent fundamental
dimensions.  Then the relationship can be written as an equation among
$n-k$ dimensionless groups $\Pi_1,\ldots,\Pi_{n-k}$:
\[
  f(\Pi_1, \Pi_2, \ldots, \Pi_{n-k}) = 0.
\]
\end{theorem}

\begin{remark}
The number $k$ is the rank of the \textbf{dimensional matrix}, whose
columns contain the exponents of the fundamental dimensions for each
quantity.
\end{remark}

\subsection{Practical Procedure}

\begin{enumerate}
  \item List all $n$ quantities involved in the problem.
  \item Build the dimensional matrix and compute its rank $k$.
  \item Form $n-k$ dimensionless groups $\Pi_j$ by combining quantities.
  \item Write the physical law as
        $\Pi_1 = \phi(\Pi_2,\ldots,\Pi_{n-k})$.
\end{enumerate}

\begin{example}[Period of a simple pendulum]
The period $\tau$ depends on the length $\ell$, mass $m$,
gravitational acceleration $g$, and amplitude $\theta_0$
(dimensionless).  The dimensioned quantities are $\tau, \ell, m, g$
($n=4$).  Fundamental dimensions: $M, L, T$ ($k=3$).  Dimensional
matrix:
\[
  \begin{pmatrix}
    & \tau & \ell & m & g \\
    M & 0 & 0 & 1 & 0 \\
    L & 0 & 1 & 0 & 1 \\
    T & 1 & 0 & 0 & -2
  \end{pmatrix}
\]
Rank $k=3$, giving $n-k=1$ dimensionless group:
\[
  \Pi_1 = \tau \sqrt{g/\ell}.
\]
Conclusion: $\tau = C(\theta_0)\sqrt{\ell/g}$ for some function $C$.
For small oscillations, $C\approx 2\pi$.  Mass does not appear.
\end{example}

% ------------------------------------------------------------
\section{Non-dimensionalisation of Equations}
\label{sec:nondim}
% ------------------------------------------------------------

\begin{definition}[Non-dimensionalisation]
To \textbf{non-dimensionalise} an equation means to introduce
dimensionless variables by dividing each quantity by a characteristic
scale:
\[
  \tilde{x} = \frac{x}{x_c}, \qquad
  \tilde{t} = \frac{t}{t_c}, \qquad
  \tilde{u} = \frac{u}{u_c}.
\]
The resulting equation involves only \textbf{dimensionless numbers}
that reveal the relative importance of different terms.
\end{definition}

\begin{example}[Damped harmonic oscillator]
The equation
\[
  m\ddot{x} + c\dot{x} + kx = 0
\]
involves mass $m$, damping $c$, stiffness $k$.  Let $x_c$ be a
characteristic amplitude and $t_c=\sqrt{m/k}$.  With $\tilde{x}=x/x_c$,
$\tilde{t}=t/t_c$:
\[
  \frac{d^2\tilde{x}}{d\tilde{t}^2}
  + 2\zeta\,\frac{d\tilde{x}}{d\tilde{t}}
  + \tilde{x} = 0,
  \qquad \zeta = \frac{c}{2\sqrt{mk}}.
\]
The only remaining parameter is the \textbf{damping ratio} $\zeta$:
\begin{itemize}
  \item $\zeta < 1$: underdamped (oscillatory decay),
  \item $\zeta = 1$: critically damped,
  \item $\zeta > 1$: overdamped (no oscillation).
\end{itemize}
\end{example}

\begin{keyformulas}[title=\textbf{Classic Dimensionless Numbers}]
\begin{center}
\begin{tabular}{lcc}
\toprule
Name & Expression & Meaning \\
\midrule
Reynolds & $\mathrm{Re}=\rho v L/\mu$ & inertia / viscosity \\
P\'eclet & $\mathrm{Pe}=vL/D$ & advection / diffusion \\
Mach & $\mathrm{Ma}=v/c_s$ & velocity / sound speed \\
Froude & $\mathrm{Fr}=v/\sqrt{gL}$ & inertia / gravity \\
Damk\"ohler & $\mathrm{Da}=\tau_{\text{transp}}/\tau_{\text{react}}$ & transport / reaction \\
\bottomrule
\end{tabular}
\end{center}
\end{keyformulas}

% ------------------------------------------------------------
\section{Scaling Laws and Similitude}
\label{sec:similitude}
% ------------------------------------------------------------

\begin{definition}[Similitude]
Two physical systems are said to be in \textbf{similitude} if they
share the same dimensionless numbers.  The study of a reduced model
(prototype) can then be transferred to the full-scale system.
\end{definition}

\begin{theorem}[Scale invariance]
If a quantity $y$ depends on variables $x_1,\ldots,x_n$ through a
power law:
\[
  y = C\,x_1^{a_1}\,x_2^{a_2}\cdots x_n^{a_n},
\]
then this law is invariant under change of units if and only if the
exponents satisfy the dimensional constraints.
\end{theorem}

\begin{example}[Drag law]
The drag force $F_D$ on a body depends on $\rho$ (fluid density),
$v$ (velocity), $L$ (characteristic length), and $\mu$ (viscosity).
We have $n=5$, $k=3$ ($M,L,T$), hence $n-k=2$ groups:
\[
  \Pi_1 = \frac{F_D}{\rho v^2 L^2}, \qquad
  \Pi_2 = \frac{\rho v L}{\mu} = \mathrm{Re}.
\]
Thus $C_D = F_D/(\tfrac{1}{2}\rho v^2 A)$ is a function of Re alone
(for a given geometry).
\end{example}

% ------------------------------------------------------------
\section{Application: Diffusion Equation}
\label{sec:diffusion_dim}
% ------------------------------------------------------------

Consider the one-dimensional diffusion equation:
\[
  \frac{\partial u}{\partial t} = D\frac{\partial^2 u}{\partial x^2},
  \qquad u(x,0) = u_0(x),
\]
on a domain of length $L$.  Set $\tilde{x}=x/L$, $\tilde{t}=Dt/L^2$,
$\tilde{u}=u/U$.  The equation becomes:
\[
  \frac{\partial\tilde{u}}{\partial\tilde{t}}
  = \frac{\partial^2\tilde{u}}{\partial\tilde{x}^2}.
\]
All parameters have disappeared: the characteristic diffusion time
is $t_c = L^2/D$.

\begin{codeblock}[title=\textbf{Diffusion time in Python}]
\begin{minted}{python}
import numpy as np
import matplotlib.pyplot as plt
from scipy.integrate import solve_ivp

L = 1.0; D = 0.01; N = 100
dx = L / (N + 1)
x = np.linspace(dx, L - dx, N)

u0 = np.exp(-((x - 0.5)**2) / (2 * 0.02**2))

def rhs(t, u):
    d2u = np.zeros_like(u)
    d2u[0] = (0 - 2*u[0] + u[1]) / dx**2
    d2u[-1] = (u[-2] - 2*u[-1] + 0) / dx**2
    d2u[1:-1] = (u[:-2] - 2*u[1:-1] + u[2:]) / dx**2
    return D * d2u

sol = solve_ivp(rhs, [0, 5], u0, t_eval=[0, 0.5, 1, 2, 5],
                method='RK45', max_step=0.01)

plt.figure(figsize=(8, 5))
for i, t_val in enumerate(sol.t):
    plt.plot(x, sol.y[:, i], label=f'$t = {t_val}$')
plt.xlabel('$x$'); plt.ylabel('$u(x,t)$')
plt.legend(); plt.title('1D Diffusion')
plt.tight_layout(); plt.savefig('diffusion_1d.pdf')
\end{minted}
\end{codeblock}

% ------------------------------------------------------------
\section{Regular Perturbations and Orders of Magnitude}
\label{sec:perturbations}
% ------------------------------------------------------------

Non-dimensionalisation often reveals a small parameter
$\varepsilon\ll 1$ enabling a \textbf{perturbation expansion}.

\begin{definition}[Regular perturbation expansion]
One seeks the solution in the form
\[
  \tilde{u} = \tilde{u}_0 + \varepsilon\,\tilde{u}_1
             + \varepsilon^2\,\tilde{u}_2 + \cdots
\]
and identifies terms order by order.
\end{definition}

\begin{example}
The weakly nonlinear oscillator:
\[
  \ddot{x} + x + \varepsilon x^3 = 0, \qquad 0<\varepsilon\ll 1.
\]
At leading order: $\ddot{x}_0 + x_0 = 0 \Rightarrow x_0 = A\cos t$.
At first order:
\[
  \ddot{x}_1 + x_1 = -x_0^3 = -A^3\cos^3 t
  = -\frac{A^3}{4}(3\cos t + \cos 3t).
\]
The $\cos t$ term is resonant (secular term), signalling that the
naive expansion breaks down at long times.  One must use multiple
scales or the Lindstedt--Poincar\'e method.
\end{example}

\begin{warning}[title=\textbf{Secular terms}]
A regular perturbation expansion may contain terms that grow without
bound in time (secular terms).  This invalidates the approximation at
long times and requires specialised techniques.
\end{warning}

% ------------------------------------------------------------
\section{Order-of-Magnitude Analysis}
% ------------------------------------------------------------

\begin{proposition}[Simplification by orders of magnitude]
After non-dimensionalisation, if a term $\varepsilon\,f(\tilde{x})$ is
multiplied by $\varepsilon\ll 1$ while the other terms are $O(1)$,
it may be neglected as a first approximation.
\end{proposition}

\begin{example}[Low-Reynolds-number flow]
For an incompressible Newtonian fluid, the non-dimensionalised
Navier--Stokes equations contain the term
$\frac{1}{\mathrm{Re}}\nabla^2\mathbf{v}$.  If $\mathrm{Re}\gg 1$,
viscosity is negligible (Euler flow).  If $\mathrm{Re}\ll 1$, inertia
is negligible (Stokes flow):
\[
  \nabla p = \mu\nabla^2\mathbf{v}.
\]
\end{example}

% ------------------------------------------------------------
\section{Exercises}
% ------------------------------------------------------------

\begin{exercise}
Using dimensional analysis, determine how the oscillation frequency
$f$ of a spherical droplet depends on its mass $m$, radius $R$, and
surface tension $\gamma$.
\end{exercise}

\begin{exercise}
Non-dimensionalise the Lotka--Volterra equations:
\[
  \dot{x} = ax - bxy, \qquad \dot{y} = -cy + dxy,
\]
by setting $\tilde{x}=bx/c$, $\tilde{y}=dy/a$, $\tilde{t}=at$.
Verify that the resulting system depends on a single parameter
$\alpha = c/a$.
\end{exercise}

\begin{exercise}
The sedimentation velocity $v_s$ of a sphere of radius $R$ and
density $\rho_s$ in a fluid of viscosity $\mu$ and density $\rho_f$
under gravity $g$ is given by Stokes' formula.
\begin{enumerate}
  \item Recover $v_s\propto R^2$ by dimensional analysis.
  \item For what values of Re is this formula valid?
\end{enumerate}
\end{exercise}

\begin{exercise}
Consider the heat equation with a source:
\[
  \rho c_p \frac{\partial T}{\partial t}
  = \kappa\frac{\partial^2 T}{\partial x^2} + Q_0 e^{-x/\ell}.
\]
\begin{enumerate}
  \item Identify characteristic scales $x_c$, $t_c$, $T_c$.
  \item Non-dimensionalise and identify the dimensionless numbers.
  \item Discuss the limiting regimes.
\end{enumerate}
\end{exercise}

\begin{exercise}
Using the Buckingham $\Pi$ theorem, show that the pressure drop
$\Delta p$ in a pipe of length $L$ and diameter $d$, for a fluid of
density $\rho$ and viscosity $\mu$ flowing at velocity $v$, takes
the form:
\[
  \frac{\Delta p}{\rho v^2} = \phi\!\left(\mathrm{Re},\,\frac{L}{d}\right).
\]
\end{exercise}
